%! AP Statistics — Unit 9, Lesson 9.3 — Guided Notes
\documentclass[10pt]{article}
\usepackage{apstats-article}
\usepackage{apstats-boxes}

\begin{document}

\pageheader{Unit 9, Lesson 9.3}{Guided Notes}
\namedateperiod

\begin{objectivebox}
By the end of this lesson, I will be able to:
\begin{itemize}
    \item[$\square$] \writeline
    \item[$\square$] \writeline
    \item[$\square$] \writeline
\end{itemize}
\end{objectivebox}

\begin{vocabbox}
\termblanklong{CI for slope}
\termblanklong{Interpretation of CI for slope}
\termblanklong{Justifying a claim with a CI}
\termblanklong{Width of CI for slope}
\end{vocabbox}

\begin{hookbox}
Computer output for a regression of exam score (points) on daily study time (hours) gives
a 95\% CI for the slope: $(0.14,\ 0.52)$.

A researcher claims the true slope $\beta = 0$ --- that studying more has no effect on exam
scores. Can you use this interval to evaluate that claim?
What if the researcher instead claimed $\beta = 0.5$?

\writelines{2}
\end{hookbox}

\begin{notesbox}{Interpreting a CI for Slope (UNC-4.AG)}

\textbf{Interpretation Template:}

``We are \blank{2cm}\% confident that the true slope of the population regression line
relating \blank{4cm} to \blank{4cm} for \blank{4cm} is between \blank{2cm} and
\blank{2cm} \blank{3cm}.''

\vspace{0.1in}
\textbf{Required elements for full AP credit:}
\begin{enumerate}
    \item \blank{6cm} (the confidence level)
    \item \blank{8cm} (both endpoints with correct units)
    \item \blank{8cm} (reference to the \emph{population}, not just the sample)
\end{enumerate}

\vspace{0.1in}
\textbf{Worked Example.} A 95\% CI for the slope relating exam score to daily study time
for high school students is $(0.14,\ 0.52)$ points per hour.

\textit{Complete interpretation:} \writelines{2}

\end{notesbox}

\begin{notesbox}{Justifying a Claim Using a CI (UNC-4.AH)}

\textbf{Logic:}
\begin{itemize}
    \item If the claimed value $\beta_0$ is \textbf{NOT} in the CI $\Rightarrow$
          \blank{7cm}
    \item If the claimed value $\beta_0$ \textbf{IS} in the CI $\Rightarrow$
          \blank{7cm}
\end{itemize}

\vspace{0.1in}
\textbf{Special case --- testing whether $\beta = 0$:}
\begin{itemize}
    \item If 0 is \textbf{not} in the CI $\Rightarrow$ \blank{6cm}
    \item If 0 \textbf{is} in the CI $\Rightarrow$ \blank{6cm}
\end{itemize}

\vspace{0.1in}
\textbf{Practice Table.} 95\% CI for slope $= (0.14,\ 0.52)$

\begin{center}
\renewcommand{\arraystretch}{1.6}
\begin{tabular}{p{3cm}p{3.5cm}p{5.5cm}}
    \textbf{Claimed value} & \textbf{In CI? (Yes/No)} & \textbf{What can we conclude?} \\
    \hline
    $\beta = 0$   & \blank{3cm} & \writeline \\
    $\beta = 0.5$ & \blank{3cm} & \writeline \\
    $\beta = 1.0$ & \blank{3cm} & \writeline \\
\end{tabular}
\end{center}

\end{notesbox}

\begin{notesbox}{Effect of Sample Size on Width (UNC-4.AI)}

\textbf{Fill in:}

``When all other things remain the same, as sample size \blank{3cm}, the standard error
$SE_b$ \blank{3cm}, and the confidence interval for slope becomes \blank{3cm}.''

\vspace{0.05in}
\textbf{Why?} \writelines{2}

\end{notesbox}

\end{document}
